% problem-000344 8 含氧芳香酸结构
\section*{8. 含氧芳香酸结构}
\textit{下列为分问。}

\subsection*{8-1}
化合物A是⼀种重要化⼯产品，⽤于⽣产染料、光电材料和治疗疣的药物等。A由第⼀、⼆周期元素组成，⽩⾊晶体，摩尔质量 $1 1 4 . 0 6 \mathrm { g } \mathrm { m o l ^ { - 1 } }$ ，熔点 $2 9 3 \mathrm { ~ } ^ { \circ } \mathrm { C }$ ，酸常数 $\mathrm { . p } K _ { a 1 } = 1 . 5 , \mathrm { p } K _ { a 2 } = 3 . 4 ;$ ；酸根离⼦ $\mathbf { A } ^ { 2 - }$ 中同种化学键是等⻓的，存在⼀根四重旋转轴。

\subsection*{8-1}
-1 画出A的结构简式。

\subsection*{8-1}
-2 为什么 $\mathbf { A } ^ { 2 - }$ 离⼦中同种化学键是等⻓的？

\subsection*{8-1}
-3 A2−离⼦有⼏个镜⾯？

\subsection*{8-2}
顺（反）丁烯⼆酸的四个酸常数为 $1 . 1 7 \times 1 0 ^ { - 2 }$ ， $9 . 3 \times 1 0 ^ { - 4 }$ $2 . 9 \times 1 0 ^ { - 5 }$ 和 $2 . 6 0 \times 1 0 ^ { - 7 }$ 。指出这些常数分别是哪个酸的⼏级酸常数，并从结构与性质的⻆度简述你做出判断的理由。

\subsection*{8-3}
氨基磺酸是⼀种重要的⽆机酸，⽤途⼴泛，全球年产量逾40万吨，分⼦量为 $9 7 . 0 9 _ { \circ }$ 晶体结构测定证实分⼦中有三种化学键，键⻓分别为 102、144 和 176 pm。氨基磺酸易溶于⽔，在⽔中的酸常数 $K _ { a } \ =$ $1 . 0 1 \times 1 0 ^ { - 1 }$ ，其酸根离⼦中也有三种键⻓，分别为 100、146 和 167 pm。

\subsection*{8-3}
-1 计算0.10 M氨基磺酸⽔溶液的 $\mathrm { p H }$ 值。

\subsection*{8-3}
-2 从结构与性质的关系解释为什么氨基磺酸是强酸？

\subsection*{8-3}
-3 氨基磺酸可⽤作游泳池消毒剂氯⽔的稳定剂，这是因为氨基磺酸与 $\mathrm { C l } _ { 2 }$ 的反应产物⼀氯代物能缓慢释放次氯酸。写出形成⼀氯代物以及⼀氯代物与⽔反应的⽅程式。
